\chapter{Coverage-Driven Test Generation}

\section{Adeguatezza dei test}

Il concetto di \textbf{adeguatezza dei test} nasce dall'esigenza di capire se e quando sia possibile smettere di testare un \textbf{System Under Test} (SUT).

Le domande principali sono:

\begin{itemize}
    \item i miei test sono sufficientemente buoni?
    \item il SUT è stato testato in modo abbastanza meticoloso?
    \item i test sono adeguati rispetto agli obiettivi prefissati?
\end{itemize}

Le tecniche di adequacy servono quindi a stimare, nel modo più oggettivo possibile, la qualità dell'attività di testing. In pratica, permettono di definire criteri di accettazione degli artefatti prodotti, cioè soglie o condizioni minime per dire che una test suite è sufficientemente adeguata.

\begin{notaprof}
L'adeguatezza viene spesso usata alla fine della catena di testing, come controllo a posteriori per dimostrare che il lavoro è stato fatto bene. In questa lezione, però, il problema viene ribaltato: l'adeguatezza può essere usata anche come \textbf{obiettivo guidato} per generare automaticamente i test fin dall'inizio.
\end{notaprof}

\begin{notaesame}
L'adequacy non dimostra l'assenza di bug. Serve a misurare quanto una test suite soddisfi certi criteri espliciti di esplorazione o verifica del SUT.
\end{notaesame}

\section{Functional vs Structural Adequacy}

Le tecniche di adequacy possono essere divise in due grandi famiglie: \textbf{functional adequacy} e \textbf{structural adequacy}.

La \textbf{functional adequacy} è un criterio di tipo black-box. Si basa su requisiti, specifiche o modelli comportamentali del sistema. L'attenzione è rivolta ai comportamenti osservabili del SUT, senza considerare direttamente il codice sorgente.

\begin{notaslide}
Il limite principale della functional adequacy è che non riesce a intercettare errori dovuti a dettagli implementativi introdotti dallo sviluppatore.
\end{notaslide}

La \textbf{structural adequacy}, invece, è un criterio di tipo white-box. Si basa sull'implementazione eseguibile del SUT e cerca di attraversare la maggior parte possibile del codice reale.

Il suo limite è duplice: da un lato non può scoprire funzionalità completamente mancanti, perché se una funzionalità non è implementata non esiste codice da coprire; dall'altro, non esiste una garanzia matematica secondo cui un'alta coverage implichi necessariamente l'esposizione di una failure.

\begin{notaprof}
Con la coverage strutturale misuriamo la nostra confidenza nell'esplorazione del sistema, ma non misuriamo direttamente il reale modo d'uso funzionale del SUT né la capacità effettiva di trovare difetti.
\end{notaprof}

\begin{notaesame}
Alta coverage non significa automaticamente buona test suite. Una riga può essere eseguita senza essere realmente verificata da un'oracolo significativo.
\end{notaesame}

\section{Generazione automatica white-box}

La \textbf{generazione automatica white-box} parte dall'idea di usare il codice del SUT come input per produrre automaticamente casi di test.

Invece di progettare manualmente test con l'obiettivo di coprire il codice, si fornisce a un algoritmo un criterio di coverage da massimizzare. Lo strumento analizza il codice e genera automaticamente test, per esempio classi JUnit.

\begin{notaslide}
Un esempio mostrato nelle slide riguarda una funzione come \texttt{areBothParamsPositive(int i, int j)}, contenente rami logici basati su condizioni come \texttt{i > 0} e \texttt{j > 0}. Uno strumento coverage-driven può generare automaticamente test JUnit per coprire i branch della funzione.
\end{notaslide}

Tuttavia, bisogna distinguere tra \textbf{generare input} e \textbf{generare veri test}. Generare input significa trovare valori che stimolano il SUT e attraversano certe parti del codice. Generare un vero test significa anche produrre un \textbf{oracle}, cioè un'asserzione che stabilisca quale comportamento sia corretto.

\begin{notaprof}
Molti strumenti automatici sono bravi a generare stimoli, cioè valori di input e sequenze di chiamate. Il problema più difficile è generare anche le \texttt{assert}, perché senza oracle il test esegue il codice ma non verifica davvero la correttezza del risultato.
\end{notaprof}

\section{Approcci evolutivi}

Gli approcci evolutivi alla generazione dei test rientrano nella \textbf{Search-Based Software Engineering}. L'idea è trattare la generazione dei test come un problema di ricerca: si parte da una popolazione di soluzioni candidate e la si evolve progressivamente verso test migliori.

Il processo generale include:

\begin{itemize}
    \item una \textbf{popolazione iniziale} di test candidati, spesso generati in modo casuale;
    \item una \textbf{fitness function}, che misura la bontà di un test rispetto all'obiettivo, per esempio code coverage o fault coverage;
    \item operazioni di \textbf{crossover}, che combinano due test candidati per produrre nuovi test;
    \item operazioni di \textbf{mutazione dei test}, che alterano localmente un test;
    \item una fase di \textbf{selezione}, in cui vengono mantenuti i candidati migliori;
    \item un criterio di arresto, come il raggiungimento dell'ottimalità o l'esaurimento del budget computazionale.
\end{itemize}

\begin{notaprof}
La mutazione dei test serve anche a evitare che l'algoritmo si blocchi in ottimi locali. Alterando casualmente una soluzione, l'algoritmo può esplorare zone dello spazio di ricerca che altrimenti non raggiungerebbe.
\end{notaprof}

Il vantaggio principale di questi approcci è la loro genericità: possono essere automatizzati e applicati a molti SUT diversi. I limiti riguardano invece la difficoltà di modellare il software come problema evolutivo, la scelta della fitness function, la taratura degli iper-parametri e, soprattutto, la generazione degli oracle.

\begin{notaesame}
Negli approcci evolutivi, ottimizzare la coverage non basta. Una test suite generata automaticamente deve anche contenere asserzioni utili e non fragili.
\end{notaesame}

\section{Definizione automatica degli oracle}

La definizione automatica degli \textbf{oracle} è uno dei problemi centrali della generazione white-box.

Trovare una sequenza di chiamate che raggiunge una certa riga di codice è utile solo in parte. Se alla fine il test non controlla l'effetto prodotto, allora non è realmente in grado di rivelare una failure.

Le asserzioni possono essere generate osservando:

\begin{itemize}
    \item valori di ritorno dei metodi;
    \item stato pubblico dell'istanza;
    \item API pubbliche del SUT;
    \item valori osservabili dopo l'esecuzione del test.
\end{itemize}

Tuttavia, questa generazione automatica può produrre problemi. Alcune asserzioni possono essere irrilevanti, superflue o troppo legate a dettagli interni dell'implementazione. In questi casi, il test diventa fragile: può fallire dopo un refactoring corretto, anche se il comportamento funzionale del sistema non è cambiato.

\begin{notaslide}
Le slide evidenziano il rischio di generare troppe asserzioni o asserzioni irrilevanti, che aumentano il rumore invece di migliorare davvero la qualità della test suite.
\end{notaslide}

\begin{notaprof}
Lo sviluppatore deve comunque controllare le asserzioni generate. Un'asserzione prodotta automaticamente non è automaticamente significativa: deve essere collegata a un requisito, a una proprietà attesa o a un comportamento realmente rilevante.
\end{notaprof}

\section{Mutazioni e asserzioni}

Per generare asserzioni più utili, alcuni approcci coverage-driven sfruttano idee derivate dal \textbf{Mutation Testing}.

\begin{notaesame}
Bisogna distinguere chiaramente tra \textbf{mutazione dei test} e \textbf{mutazione del SUT}. La mutazione dei test riguarda modifiche casuali alla popolazione di test candidati. La mutazione del SUT riguarda invece alterazioni artificiali del codice sorgente per simulare fault.
\end{notaesame}

L'idea è assumere che il SUT originale sia corretto, introdurre artificialmente una mutazione nel SUT e poi eseguire la stessa sequenza di test sia sul programma originale sia sul mutante.

Se il comportamento osservabile diverge, il tool può usare quella divergenza per generare un'asserzione. In altre parole, l'asserzione viene costruita in modo da accettare il comportamento del SUT originale e rifiutare quello del mutante.

Il processo può essere sintetizzato così:

\begin{enumerate}
    \item si genera una sequenza di test;
    \item si esegue la sequenza sul SUT originale;
    \item si esegue la stessa sequenza su un mutante del SUT;
    \item si confrontano ritorni, stato e valori osservabili;
    \item si genera un'asserzione che cattura la differenza significativa.
\end{enumerate}

\begin{notaprof}
Un mutante viene rilevato solo se esiste un oracle capace di osservare la divergenza. Se il test attraversa il codice mutato ma non controlla il valore modificato, il mutante può sopravvivere.
\end{notaprof}

\section{Test come sequenze di istruzioni}

Negli approcci evolutivi, un test viene modellato come una \textbf{sequenza di istruzioni}. Questa modellazione permette agli algoritmi di generare, combinare e modificare test in modo automatico.

\begin{notaslide}
Le slide distinguono tra \textbf{constructor statements}, \textbf{method statements}, \textbf{field statements} e \textbf{primitive statements}.
\end{notaslide}

I \textbf{constructor statements} creano nuove istanze, per esempio tramite \texttt{new}. I \textbf{method statements} invocano metodi su oggetti già creati. I \textbf{field statements} accedono a campi pubblici, mentre i \textbf{primitive statements} definiscono valori primitivi, come interi, booleani o stringhe.

La generazione deve rispettare alcune regole. Ogni istruzione produce una nuova variabile riutilizzabile. I parametri passati a un metodo devono riferirsi a variabili già dichiarate, altrimenti il codice generato sarebbe sintatticamente non valido.

Inoltre, possono essere coinvolte anche classi dipendenti dal SUT. Questo è necessario perché molte API richiedono oggetti complessi come parametri; quindi, prima di invocare il metodo da testare, bisogna costruire uno stato valido e coerente.

\begin{notaesame}
Generare test automatici non significa solo scegliere valori primitivi. Spesso bisogna costruire intere sequenze di oggetti e chiamate per portare il SUT in uno stato utile al test.
\end{notaesame}

\section{EvoSuite}

\textbf{EvoSuite} è uno strumento per la generazione automatica di test suite JUnit. Nasce da ricerche accademiche e implementa un approccio coverage-based alla generazione dei test.

\begin{notaslide}
EvoSuite lavora a livello di \textbf{bytecode Java}, cioè sui file \texttt{.class}. Per questo non richiede necessariamente il codice sorgente del SUT o delle librerie.
\end{notaslide}

Il suo obiettivo è generare automaticamente una test suite cercando di:

\begin{itemize}
    \item massimizzare la code coverage;
    \item generare test JUnit eseguibili;
    \item produrre asserzioni;
    \item minimizzare il numero di asserzioni non necessarie.
\end{itemize}

EvoSuite può essere usato in diverse modalità operative: stand-alone da riga di comando, tramite plugin per Eclipse, tramite plugin per IntelliJ o tramite plugin Maven.

\begin{notaesame}
EvoSuite mostra bene il compromesso della generazione automatica: può produrre molti test e aumentare la coverage, ma le asserzioni generate devono essere controllate criticamente.
\end{notaesame}

\section{Collegamento con il progetto}

Nel progetto è utile distinguere tra test scritti manualmente, test generati con strumenti randomici come Randoop, test prodotti da LLM e test generati da strumenti coverage-driven come EvoSuite.

I test manuali permettono di ragionare sulle specifiche e sulle classi di equivalenza. Randoop genera sequenze random guidate dal feedback. Gli LLM possono produrre test a partire da codice e prompt, ma richiedono validazione. EvoSuite e gli strumenti coverage-driven, invece, usano la coverage come obiettivo esplicito di generazione.

\begin{notaesame}
Il peggior errore è fidarsi solo della coverage. Una test suite automatica può coprire molto codice ma contenere oracle deboli, asserzioni fragili o test poco significativi.
\end{notaesame}

Nel report di progetto bisogna documentare:

\begin{itemize}
    \item il tool usato;
    \item la configurazione adottata;
    \item il criterio di coverage considerato;
    \item i test generati;
    \item il tipo e la qualità delle asserzioni;
    \item i limiti osservati;
    \item eventuali test scartati o modificati.
\end{itemize}

\begin{notaprof}
Includere test automatici che fanno molte chiamate ma non contengono nessuna \texttt{assert} significativa significa non aver compreso lo scopo del testing. Un test deve stimolare il SUT, ma deve anche verificare il risultato.
\end{notaprof}

